% !PRINTED 10
\weberpage{10}
%<<CONT
$A''$, namely fractions of the forms $a+a'$ and $b+b'$. Further,
$\alpha''$ is greater than $\alpha$ and than $\alpha'$. For $A''$
first contains all of $A$, and if $a'$ is an arbitrary fraction in
$A'$, one can choose $a$ in $A$ so that the element $a+a'$ of $A''$ is
contained in $B$; hence $A''$ is more comprehensive than $A$. The cuts
produced by rational fractions $\mu,\mu'$ give, by addition, the cut
produced by $\mu+\mu'$.

\begin{center}
\rule{0.18\linewidth}{0.4pt}
\end{center}

We now pass to the definition of \gesperrt{ratios}, which since
antiquity have been regarded as the foundation of the theory of
numbers, and here follow Euclid first of all.
\setcounter{footnote}{0}\footnote{Elements, Book V.}

If one combines the elements of a measurable set $\mathfrak M$ into
pairs, and regards these pairs themselves as elements, then a new
manifold arises. We denote such a pair by $a:b$, or also by $a/b$; but,
when $a$ and $b$ are different elements, we distinguish $a:b$ from
$b:a$, and call $a$ the \gesperrt{numerator} and $b$ the
\gesperrt{denominator} of $a:b$. We shall call these pairs ratios, and
shall now order and make measurable this new set.

Assume first that two whole numbers $m,n$ exist such that $na=mb$, as
is always the case, for example, when $\mathfrak M$ is the system of
natural numbers, or when $a,b$ are two commensurable segments. Then,
if $p,q$ are two other whole numbers, $qa=pb$ holds if and only if
$mq=np$. The pair of numbers $p,q$ is completely determined by this
requirement if the condition is added that $p,q$ are to be relatively
prime. Then, if $h$ is an arbitrary whole number, $m=hp$, $n=hq$. In
this case we call the ratio $a:b$ \gesperrt{rational} and set it equal
to the rational fraction $m:n$ or $p:q$. Accordingly these rational
fractions may be regarded as ratios of whole numbers.

All rational ratios equal to one another form a \gesperrt{rational
number}, and the rational numbers form, like the rational fractions,
an ordered, dense, and measurable manifold.
